\chapter{The Non-Perturbative Decoupling Theorem}
\label{ch:fix8_decoupling}

\section{Context}
The Adjoint Interpolation strategy relies on the fact that as the fermion mass $M \to \infty$, the theory converges to Pure Yang-Mills. We must rigorously prove that the heavy fermions decouple without leaving non-local "scars" or shifting the vacuum parameters, upgrading the perturbative Appelquist-Carazzone theorem to a non-perturbative operator convergence.

\section{Theorem N.1 (Norm Resolvent Convergence)}
Let $H(M)$ be the Hamiltonian of Adjoint QCD with mass $M$, and $H_{YM}$ be the Pure Yang-Mills Hamiltonian. As $M \to \infty$:
\begin{equation}
\| (H(M) - z)^{-1} - (H_{YM} - z)^{-1} \| \le \frac{C}{M}
\end{equation}
This implies uniform convergence of the spectrum, ensuring the mass gap $\Delta(M)$ does not collapse discontinuously at $M=\infty$.

\section{Proof Derivation}
\subsection{1. Feshbach-Schur Map}
We decompose the Hilbert space $\mathcal{H} = P\mathcal{H} \oplus Q\mathcal{H}$, where $P$ projects onto the pure gauge sector (fermion number 0) and $Q$ onto states with fermions. The effective Hamiltonian on the gauge sector is:
\begin{equation}
H_{eff} = H_{YM} - P V Q (Q H(M) Q - z)^{-1} Q V P
\end{equation}.

\subsection{2. Kinetic Gap}
The operator $Q H(M) Q$ contains the fermion mass term, so $Q H(M) Q \ge M$. For energies $z \ll M$, the resolvent $(Q H Q - z)^{-1}$ can be expanded in a Neumann series (hopping parameter expansion).

\subsection{3. Interaction Bound}
The interaction term $V$ represents pair creation/annihilation. Its norm is bounded on the lattice. The shift in the effective Hamiltonian is bounded by:
\begin{equation}
\| \Delta H \| \le \frac{\|V\|^2}{M - z}
\end{equation}.

\subsection{4. Spectral Stability}
Since the operators converge in norm, the spectral gap $\Delta(M)$ converges to $\Delta_{YM}$. Since $\Delta(M) > 0$ for all finite $M$ (via the interpolation argument), the limit $\Delta_{YM}$ is strictly positive.



